
\subsubsection{3.1 Overview}

Our methodology consists of a three-phase pipeline applied to a random sample of Python files:

\begin{itemize}
\item \textbf{Phase A (Pattern):} Does the file contain any \texttt{>>>} substring?
\item \textbf{Phase B (AST):} Does the file contain AST-parseable doctest examples?
\item \textbf{Phase C (Subprocess):} Do the doctest examples pass subprocess execution in a clean environment?
\end{itemize}

\subsubsection{3.2 Data Source and Sampling}

\textbf{Dataset:} \texttt{codeparrot/codeparrot-clean-valid}, a curated Python-only corpus publicly
accessible on HuggingFace Hub. Used as a proxy for \texttt{bigcode/the-stack-dedup}, which requires
HuggingFace Hub access approval and was unavailable in the execution environment. Both datasets
use the same \texttt{content} field schema and contain Python-only code. The fallback is confirmed in
the experiment log (\texttt{experiment.log}): ''the-stack-dedup unavailable (gated). Falling back to
codeparrot/codeparrot-clean-valid.''

\textbf{Sample:} 10,000 randomly sampled Python files (streaming + first-N from shuffled stream,
\texttt{seed=42}), consistent with The Stack paper's own analysis methodology \citep{Kocetkov2022}.

\textbf{Quality Pre-filter:} The Stack paper quality filters applied prior to sampling: average
line length $\leq$ 100 characters, maximum line length $\leq$ 1,000 characters, alphanumeric fraction
$\geq$ 0.25.

\subsubsection{3.3 Three-Phase Filtering Pipeline}

\textbf{Phase A: Pattern Detection.} \texttt{''>>>'' in source\textbackslash \_code} --- O(n) substring check providing
an upper bound on doctest-bearing files. Files without \texttt{>>>} are immediately excluded.

\textbf{Phase B: AST Extraction.}

\begin{verbatim}
# Inside check_phase_b(source_code: str) -> bool
tree = ast.parse(source_code)
parser = doctest.DocTestParser()
for node in ast.walk(tree):
    if isinstance(node, (ast.FunctionDef, ast.AsyncFunctionDef, ast.ClassDef, ast.Module)):
        docstring = ast.get_docstring(node)
        if docstring and ">>>" in docstring:
            examples = parser.get_examples(docstring)
            if examples:
                return True  # Phase B positive
\end{verbatim}

Phase B eliminates files with syntax errors and non-doctest \texttt{>>>} occurrences (e.g., shell
prompts in comments), filtering approximately 35\% of Phase A positives before the expensive
subprocess step.

\textbf{Phase C: Subprocess Execution.} Source code is base64-encoded and executed in an isolated
subprocess with a 5-second timeout:

\begin{verbatim}
# base64 encoding eliminates shell quoting errors across all 10,000 files
encoded = base64.b64encode(source.encode("utf-8")).decode("ascii")
wrapper = (
    "import base64, types, doctest, sys\n"
    f"src = base64.b64decode('{encoded}').decode('utf-8')\n"
    "mod = types.ModuleType('__scanned__')\n"
    "exec(compile(src, '__scanned__', 'exec'), mod.__dict__)\n"
    "results = doctest.testmod(mod, verbose=False)\n"
    "sys.exit(0 if results.failed == 0 else 1)\n"
)
result = subprocess.run([sys.executable, "-c", wrapper],
                        capture_output=True, timeout=5)
\end{verbatim}

\textbf{Parallelism:} \texttt{ProcessPoolExecutor} with 4 workers for Phase C. Only Phase B positive
files are submitted, limiting subprocess overhead.

\textbf{Error Classification:} Failed Phase C files are classified by error type: \texttt{import\textbackslash \_error}
(ImportError or ModuleNotFoundError in stderr), \texttt{assertion\textbackslash \_error} (AssertionError or
''Failed example''), \texttt{wrong\textbackslash \_output} (DocTestFailure or UnexpectedException), or \texttt{exception}
(all other non-zero exit codes).

\subsubsection{3.4 Gate Evaluation}

Gate thresholds determine the downstream SFT experiment design:

\begin{table}[htbp]
\centering
\resizebox{\columnwidth}{!}{%
\begin{tabular}{lll}
\toprule
Result & `doctest\_executable\_rate` & Decision \\
\midrule
**PASS** & $\geq$ 3.0\% & 3-condition H-E1 (unfiltered / compile-only / doctest) \\
**SCOPE** & 1.0--3.0\% & Reduced token budget \\
**PIVOT** & < 1.0\% & 2-condition H-E1 (unfiltered / compile-only) \\
\bottomrule
\end{tabular}
}
\end{table}

The gate type is SHOULD\_WORK: gate failure does not block the pipeline but instead scopes
the downstream SFT experiment design.

\subsubsection{3.5 Planned SFT Comparison (H-E1, Pending)}

Based on the H-C1 PIVOT result, the SFT experiment is designed with two conditions and
pending execution:

\begin{itemize}
\item \textbf{Model:} Qwen2.5-Coder-1.5B \citep{Hui2024}
\item \textbf{Conditions:} Unfiltered random subsample vs. compile()-filtered subsample, equal 500M-token budget
\item \textbf{Training:} AdamW, lr=2e-5, batch=32, 3 epochs, seed=42
\item \textbf{Evaluation:} HumanEval pass@1 and MBPP pass@1 (greedy decoding, lm-evaluation-harness \citep{EvalHarness2021})
\end{itemize}

